\documentclass{article}
\usepackage[margin=28mm]{geometry}
\usepackage{amsmath,amssymb}
\usepackage{xurl}
\usepackage[hidelinks]{hyperref}
% Shared commands for notes in docs/paper-gaps/.

\DeclareUnicodeCharacter{2080}{\ifmmode{}_{0}\else\textsubscript{0}\fi}
\DeclareUnicodeCharacter{2081}{\ifmmode{}_{1}\else\textsubscript{1}\fi}
\DeclareUnicodeCharacter{2082}{\ifmmode{}_{2}\else\textsubscript{2}\fi}
\DeclareUnicodeCharacter{2099}{\ifmmode{}_{n}\else\textsubscript{n}\fi}

\newcommand{\Lean}{\textsc{Lean}}
\ExplSyntaxOn
\NewDocumentCommand{\leanid}{m}{
  \ifmmode
    \textsf{\detokenize{#1}}
  \else
    \group_begin:
    \urlstyle{sf}
    \tl_set:Nn \l_tmpa_tl {#1}
    \tl_replace_all:Nnn \l_tmpa_tl {\_} {_}
    \exp_args:NV \path \l_tmpa_tl
    \group_end:
  \fi
}
\ExplSyntaxOff
\providecommand{\tr}{\operatorname{tr}}
\newcommand{\tnleanIssueBase}{https://github.com/LionSR/TNLean/issues/}
\newcommand{\tnleanPrBase}{https://github.com/LionSR/TNLean/pull/}
\newcommand{\tnleanDocBase}{https://sirui-lu.com/TNLean/docs/}
\newcommand{\ghissue}[1]{\href{\tnleanIssueBase#1}{GitHub issue \##1}}
\newcommand{\ghpr}[1]{\href{\tnleanPrBase#1}{PR \##1}}
\newcommand{\doclink}[1]{\href{\tnleanDocBase#1.html}{\path{#1}}}

% Dirac notation and the identity operator, as in blueprint/src/macros/common.tex.
% \providecommand keeps note-local definitions authoritative.
\usepackage{dsfont}
\providecommand{\ket}[1]{|#1\rangle}
\providecommand{\bra}[1]{\langle #1|}
\providecommand{\braket}[2]{\langle #1 | #2 \rangle}
\providecommand{\ketbra}[2]{|#1\rangle\!\langle #2|}
\providecommand{\Id}{\mathds{1}}
\providecommand{\one}{\mathds{1}}
\providecommand{\C}{\mathbb{C}}
\providecommand{\MN}[1]{M_{#1}(\C)}
\providecommand{\MD}{M_D(\C)}

% Verdict marker: \gapnote{<kind>}{<status>}. Kind is the policy
% classification (clarification, local-correction, scope-restriction,
% unfaithful, false-source, open-gap); status is open, wip (elimination
% underway), resolved, or historical. Machine-read by the site index;
% prints a verdict line.
\newcommand{\gapnote}[2]{\par\noindent\textbf{Verdict:} #1 (#2).\par}

\providecommand{\gapnote}[2]{\par\noindent\textbf{Verdict:} #1 (#2).\par}
\title{The Example of Non-Commuting Low-Energy Terms in Hastings' Appendix~B}
\author{}
\date{}
\begin{document}
\maketitle
\gapnote{local-correction}{resolved}

\section{The source example}
Hastings, \emph{Solving gapped Hamiltonians locally}~\cite{Hastings2006Solving},
\href{https://arxiv.org/abs/cond-mat/0508554v3}{arXiv:cond-mat/0508554v3},
source \path{proj3new.tex}, lines 1655--1676, shows that without translation
invariance the local terms need not commute in the low-energy sector. On a
chain of spins $1/2$ it takes
\begin{align}
    H_i=S^z_i\quad(i\neq0,1,2),\qquad
    H_0=S^x_0+S^y_0,\qquad
    H_1=-S^x_0-S^z_1,\qquad
    H_2=-S^y_0+S^z_2,
    \label{eq:hastings05_low_energy_example_terms}
\end{align}
and states that this ``is a complicated way of writing the Hamiltonian
$H=\sum_{i\neq0}S^z_i$'' (line 1660), whose ground space is spanned by the two
states with every spin down except spin~$0$. It then states that $H_0$, $H_1$
and $H_2$ pairwise fail to commute in the low-energy sector, ``still $H_0$
commutes in the low energy sector with $H_i$ for $i\geq2$'' (line 1670).

\section{The two corrections}
The terms acting on site~$0$ cancel, but the sum of the terms in
(\ref{eq:hastings05_low_energy_example_terms}) is
\begin{align}
    H_0+H_1+H_2=-S^z_1+S^z_2,
    \qquad
    H=-S^z_1+\sum_{i\geq2}S^z_i,
    \label{eq:hastings05_low_energy_example_sum}
\end{align}
not $\sum_{i\neq0}S^z_i$: the sign of $S^z_1$ is reversed. The ground space of
(\ref{eq:hastings05_low_energy_example_sum}) is still two-dimensional,
spanned by the states with spin~$1$ up, every spin $i\geq2$ down and spin~$0$
free. On it $S^z_k$ is a scalar for every $k\neq0$, so the compressions of
$H_0$, $H_1$ and $H_2$ are, up to scalars, those of $S^x_0+S^y_0$, $-S^x_0$ and
$-S^y_0$, which pairwise fail to commute because $[S^x_0,S^y_0]\neq0$. The
conclusions of the example are therefore unchanged.

The range ``$i\geq2$'' at line 1670 contradicts the preceding sentence, which
lists $H_0$ and $H_2$ as non-commuting. The compression of $H_0$ commutes with
that of $H_i$ exactly for $i\geq3$, where $H_i=S^z_i$ is a scalar on the ground
space.

\section{The truncated filtered terms}
The source's claim is about the filtered local approximations: ``the $M_i$ need
not commute in the low energy sector'' (line 1656). It reaches this from the
terms $H_i$ through the remark that ``for large enough $q,l_{proj}$, the
commutator of the $M_i$ in the low energy sector is close to that of the $H_i$
in the low energy sector'' (lines 1665--1667), which is stated without
argument. In this example the passage is exact. The Hamiltonian in
(\ref{eq:hastings05_low_energy_example_sum}) acts only through the commuting
operators $S^z_i$ with $i\ne0$. Each $H_i$ is a sum of an operator on spin~$0$
and such diagonal operators, so $[H,H_i]=0$. Thus
$e^{itH}H_ie^{-itH}=H_i$, and the normalized Gaussian average is $H_i$ for
every variance.

For the chain $0,\ldots,m+2$, all terms are within distance $m+2$ of any
site. Take the interaction range to be $R=2$ and
$l_{\mathrm{proj}}\ge m+4$. The generator of the source's truncated dynamics,
$H_{\mathrm{loc}}=\sum_{|i-j|\le l_{\mathrm{proj}}-2}H_j$, is then exactly $H$.
Consequently $M_i=H_i$ for every $i$, and all the noncommutation and
commutation conclusions hold for the actual truncated terms. This gives an
explicit sufficient radius for the source's ``large enough'' assertion;
no limiting argument in the variance is needed.

The source also leaves the chain unspecified beyond site~$0$; sites to the
left of $0$ with $H_i=S^z_i$ would only add scalars on the ground space,
and the formalized chain starts at site~$0$.

\section{The formalization}
Both corrections are local: the example is formalized with the sum
(\ref{eq:hastings05_low_energy_example_sum}) and the range $i\geq3$, on a chain
of $m+3$ sites.\footnote{\leanid{TNLean.Lattice.LowEnergyExample.hamiltonian_eq_diagonal},
    \leanid{TNLean.Lattice.LowEnergyExample.not_commute_compress_term_zero_one},
    \leanid{TNLean.Lattice.LowEnergyExample.commute_compress_term_zero_of_three_le}
    in \doclink{TNLean/Lattice/LowEnergyCommutatorExample}.}
The equalities for the filters are
\leanid{TNLean.Lattice.LowEnergyExample.gaussianFilteredTerm_term_eq} and
\leanid{TNLean.Lattice.LowEnergyExample.truncatedTerm_eq}; the conclusions for
the truncated terms are
\leanid{TNLean.Lattice.LowEnergyExample.not_commute_compress_truncatedTerm} and
\leanid{TNLean.Lattice.LowEnergyExample.commute_compress_truncatedTerm_zero_of_three_le}.

\bibliographystyle{alpha}
\bibliography{references}

\end{document}
